\documentclass[10pt,a4paper]{article}
\usepackage{amsmath,amssymb,amsthm,graphicx}
% Card-local minimal preamble (journal unit sigma-2025-030-sign-convention-for-a-infty-operations-in-bo, P2-ahl16 precedent).
%% Verbatim from the pinned source: the author theorem environments and macro block,
%% so every macro and environment of the counted spans is available.
% Omitted and why: the documentclass line (the journal class sigma is not on the host;
%% the assembler supplies the class) and page-geometry/font blocks (they fight the
%% forced xelatex engine).
\usepackage{amsmath,amssymb,amsthm,mathtools}

\numberwithin{equation}{section}
\newtheorem{thm}{Theorem}[section]
\newtheorem{cor}[thm]{Corollary}
\newtheorem{lem}[thm]{Lemma}
\newtheorem{sublem}[thm]{Sublemma}
\newtheorem{subsublemma}[thm]{Subsublemma}
\newtheorem{prop}[thm]{Proposition}
\newtheorem{ax}[thm]{Axiom}
\newtheorem{conj}[thm]{Conjecture}
\newtheorem{clm}[thm]{Claim}
\newtheorem{lemdef}[thm]{Lemma-Definition}
\newtheorem{observe}[thm]{Observation}
\theoremstyle{definition}
\newtheorem{defn}[thm]{Definition}
\newtheorem{rem}[thm]{Remark}
\newtheorem{ques}[thm]{Question}
\newtheorem{cons}[thm]{\rm\bfseries{Construction}}
\newtheorem{exm}[thm]{Example}
\newtheorem{conds}[thm]{Condition}
\newtheorem{conven}[thm]{Convention}
\newtheorem{proper}[thm]{Property}
\newtheorem{defnlem}[thm]{Definition-Lemma}
\newtheorem{assump}[thm]{Assumption}
\newtheorem{warn}[thm]{Warning}
\newtheorem{shitu}[thm]{Situation}
\newtheorem{choi}[thm]{Choice}
\newtheorem{notation}[thm]{Notation}
\newtheorem{problem}[thm]{Problem}
\newcommand{\del}{\partial}
\newcommand{\bMor}{\bf Mor}
 \newcommand{\ba}{\begin{eqnarray}}
  \newcommand{\na}{\end{eqnarray}}
 \newcommand{\disp}{\displaystyle}
\newcommand{\Mor}{{\text{Mor}}}
 \newcommand{\IMor}{{\text{IMor}}}
\newcommand{\rank}{\operatorname{rank}}
\newcommand{\IR}{\mbox{I \hspace{-0.2cm}R}}
\newcommand{\IN}{\mbox{I \hspace{-0.2cm}N}}
\renewcommand{\a}{\alpha}
 \renewcommand{\b}{\beta}
 \renewcommand{\c}{\gamma}
 \newcommand{\e}{\epsilon}
 \renewcommand{\d}{\delta}
  \newcommand{\p}{\partial}
  \newcommand{\z}{\mathfrak z}
\def\EE{{\mathcal E}}
\def\KK{{\mathcal K}}
\def \mc{\mathcal}
\def \wt{\widetilde}
\def\E{\ifmmode{\mathbb E}\else{$\mathbb E$}\fi} %
\def\N{\ifmmode{\mathbb N}\else{$\mathbb N$}\fi} %natural numbers%
\def\R{\ifmmode{\mathbb R}\else{$\mathbb R$}\fi} %real numbers
\def\Q{\ifmmode{\mathbb Q}\else{$\mathbb Q$}\fi} %rational numbers
\def\C{\ifmmode{\mathbb C}\else{$\mathbb C$}\fi} %complex numbers
\def\H{\ifmmode{\mathbb H}\else{$\mathbb H$}\fi} %complex numbers
\def\Z{\ifmmode{\mathbb Z}\else{$\mathbb Z$}\fi} %integers
\def\P{\ifmmode{\mathbb P}\else{$\mathbb P$}\fi} %
\def\T{\ifmmode{\mathbb T}\else{$\mathbb T$}\fi} %
\def\SS{\ifmmode{\mathbb S}\else{$\mathbb S$}\fi} %
\def\DD{\ifmmode{\mathbb D}\else{$\mathbb D$}\fi} %
\def\K{\ifmmode{\mathbb K}\else{$\mathbb K$}\fi}
\def\F{\ifmmode{\mathbb F}\else{$\mathbb F$}\fi} %finite field
\newcommand{\arc}[1]{%
}
\usepackage{cleveref}
\crefname{originaltheorem}{original theorem}{original theorems}
\Crefname{originaltheorem}{Original theorem}{Original theorems}
\makeatletter
\newlabel{bdryori}{{4.1}{source}{}{sourcepointer-1}{}}
\newlabel{bdryori@cref}{{[originaltheorem][0][]4.1}{[1][0][]source}{}{}{}}
\makeatother
\crefname{originaltheorem}{original theorem}{original theorems}
\Crefname{originaltheorem}{Original theorem}{Original theorems}
\makeatletter
\newlabel{corrStokes}{{2.5}{source}{}{sourcepointer-2}{}}
\newlabel{corrStokes@cref}{{[originaltheorem][0][]2.5}{[1][0][]source}{}{}{}}
\makeatother
\crefname{originalequation}{original equation}{original equations}
\Crefname{originalequation}{Original equation}{Original equations}
\makeatletter
\newlabel{rel1}{{4.1}{source}{}{sourcepointer-3}{}}
\newlabel{rel1@cref}{{[originalequation][0][]4.1}{[1][0][]source}{}{}{}}
\makeatother
\crefname{originalequation}{original equation}{original equations}
\Crefname{originalequation}{Original equation}{Original equations}
\makeatletter
\newlabel{rel2}{{4.2}{source}{}{sourcepointer-4}{}}
\newlabel{rel2@cref}{{[originalequation][0][]4.2}{[1][0][]source}{}{}{}}
\makeatother
\crefname{originaldefinition}{original definition}{original definitions}
\Crefname{originaldefinition}{Original definition}{Original definitions}
\makeatletter
\newlabel{def}{{3.3}{source}{}{sourcepointer-5}{}}
\newlabel{def@cref}{{[originaldefinition][0][]3.3}{[1][0][]source}{}{}{}}
\makeatother
\crefname{originalequation}{original equation}{original equations}
\Crefname{originalequation}{Original equation}{Original equations}
\makeatletter
\newlabel{lem}{{4.5}{source}{}{sourcepointer-6}{}}
\newlabel{lem@cref}{{[originalequation][0][]4.5}{[1][0][]source}{}{}{}}
\makeatother
\title{The Bott-Morse A-infinity operations in the de Rham model satisfy the filtered A-infinity relation}
\author{Ono, Kaoru}
\date{doi:10.3842/SIGMA.2025.030; CC BY 4.0}
\begin{document}
\maketitle
\setcounter{section}{4}
\setcounter{thm}{3}
\setcounter{equation}{7}
\noindent\textit{Editorial note.} All delivered mathematical text is the author's own, copied from the pinned journal source file sources/sigma-2025-030-sign-convention-for-a-infty-operations-in-bo/pinned-original-source.tex (Sign Convention for A-infinity-Operations in Bott-Morse Case); no statement or proof step was rewritten, reordered or removed.\par\medskip
\noindent\textit{Editorial note.} The counted proof is the author's complete argument: Named-input pointers declared as external references: Proposition bdryori (4.1, source line 564, used at line 962), Proposition corrStokes (2.5, line 241, used at line 977), equation (4.1) rel1 (line 528) and equation (4.2) rel2 (line 534), both used at line 985, Definition def (3.3, line 475, used at line 955), equation (4.5) lem (line 850, inside Lemma lemma 4.3, used at line 965); all are in-source proven results or definitions referenced by name (navigation pointers). bdryori is an in-source proven proposition; its own proof (source lines 579-782) delegates one base case to [FOOO09II, p.699] - disclosed, the reference here is to the proposition as a named result. In-span labels render from the card's own counters: Claim rel3 (line 964, renders 4.5), equation (4.8) Stokes (line 981). The statement's defining relation is an unnumbered equation* carried verbatim. Census printed field is not used; the card face renders Theorem 4.4 by source counter stepping (section 4, fourth shared-counter environment: 4.1 Proposition bdryori line 564, 4.2 Lemma compformula line 787, 4.3 Lemma lemma line 849; the 26 custom theorem environments share the [thm] counter numbered within section per the source preamble).\par\medskip
\noindent\textit{Editorial note.} The label conclusion names the counted statement, and the proof environment immediately following it in the source is delivered verbatim as its proof.\par\medskip
\noindent\textit{Editorial note.} The counted statement is carried by the author's own thm environment in the article source.\par\medskip
\noindent\textit{Editorial note.} Journal-unit assembly: card-local preamble (author theorem environments and macros verbatim), article class, no bibliography in the counted span (the span carries no \textbackslash{}cite commands).\par\medskip
\noindent\textit{Editorial scope.} the counted statement and its own proof environment only\par\medskip
\noindent\textit{Allowed standard prerequisites (editorial).}\begin{itemize}
\item Definition def (Definition 3.3, line 475) of the operations m\_k, the Bott-Morse A-infinity-operation in the de Rham model.
\item The boundary-orientation formula of Proposition bdryori (Proposition 4.1, line 564), valid under the tree-like K-system and CF-perturbation framework (cited [FOOO20] at line 514, outside the span).
\item The decomposed (1,0)-case relations: equation (4.1) rel1 (line 528), m\_\{1,0\} circ m\_\{k,B\} = (-1)\textasciicircum{}\{epsilon\} d (ev\_0\textasciicircum{}B)\_! (ev\_1\textasciicircum{}\{B*\} xi\_1 wedge ... wedge xi\_k), and equation (4.2) rel2 (line 534), m\_\{k,B\} circ m-hat\_\{1,0\} = (-1)\textasciicircum{}\{epsilon+1\} (ev\_0\textasciicircum{}B)\_! d (ev\_1\textasciicircum{}\{B*\} xi\_1 wedge ... wedge xi\_k).
\item The Stokes-type formula of Proposition corrStokes (Proposition 2.5, line 241, whose statement cites [FOOO20, Proposition 27.2] outside the span).
\item The composition formula of Lemma lemma (Lemma 4.3, line 849) with its sign kappa' (equation (4.5)), for the summands with both factors different from (1,0).
\item Standing: the graded-coderivation extension m-hat\_\{k,B\}, the shifted degree |.|', and epsilon(xi\_1,...,xi\_k).
\end{itemize}\medskip
\ifdefined\hypertarget\hypertarget{sourcepointer-1}{}\fi
\noindent\textit{Source reference.} Original theorem 4.1: operations/codex-takeover-01a12558/corpus-repair-20261010/research-line/harness-ready-20261011/sources/sigma-2025-030-sign-convention-for-a-infty-operations-in-bo/pinned-original-source.tex, source lines 564--564. Navigation pointer invoked at line(s) 962: the counted text invokes that named theorem (bdryori) of the same article as an input, which is not retained in this card, so the reference is declared as a pointer carrying the number the article itself prints for it (the section and the place of the statement inside that section, derived by the census v3.5 rule for the shared section-scoped theorem counter); the cited theorem is an allowed input and none of it is reproduced inside the span.\par\medskip
\ifdefined\hypertarget\hypertarget{sourcepointer-2}{}\fi
\noindent\textit{Source reference.} Original theorem 2.5: operations/codex-takeover-01a12558/corpus-repair-20261010/research-line/harness-ready-20261011/sources/sigma-2025-030-sign-convention-for-a-infty-operations-in-bo/pinned-original-source.tex, source lines 241--241. Navigation pointer invoked at line(s) 977: the counted text invokes that named theorem (corrStokes) of the same article as an input, which is not retained in this card, so the reference is declared as a pointer carrying the number the article itself prints for it (the section and the place of the statement inside that section, derived by the census v3.5 rule for the shared section-scoped theorem counter); the cited theorem is an allowed input and none of it is reproduced inside the span.\par\medskip
\ifdefined\hypertarget\hypertarget{sourcepointer-3}{}\fi
\noindent\textit{Source reference.} Original equation 4.1: operations/codex-takeover-01a12558/corpus-repair-20261010/research-line/harness-ready-20261011/sources/sigma-2025-030-sign-convention-for-a-infty-operations-in-bo/pinned-original-source.tex, source lines 528--528. Navigation pointer invoked at line(s) 985: the counted text cites that numbered equation of the same article, which is not retained in this card, so the reference is declared as a pointer carrying the number the article itself prints for it (the section and the place of the equation inside that section, derived by the census v3.5 rule for the section-scoped equation counter); the cited equation is an allowed input and none of it is reproduced inside the span.\par\medskip
\ifdefined\hypertarget\hypertarget{sourcepointer-4}{}\fi
\noindent\textit{Source reference.} Original equation 4.2: operations/codex-takeover-01a12558/corpus-repair-20261010/research-line/harness-ready-20261011/sources/sigma-2025-030-sign-convention-for-a-infty-operations-in-bo/pinned-original-source.tex, source lines 534--534. Navigation pointer invoked at line(s) 985: the counted text cites that numbered equation of the same article, which is not retained in this card, so the reference is declared as a pointer carrying the number the article itself prints for it (the section and the place of the equation inside that section, derived by the census v3.5 rule for the section-scoped equation counter); the cited equation is an allowed input and none of it is reproduced inside the span.\par\medskip
\ifdefined\hypertarget\hypertarget{sourcepointer-5}{}\fi
\noindent\textit{Source reference.} Original definition 3.3: operations/codex-takeover-01a12558/corpus-repair-20261010/research-line/harness-ready-20261011/sources/sigma-2025-030-sign-convention-for-a-infty-operations-in-bo/pinned-original-source.tex, source lines 475--475. Navigation pointer invoked at line(s) 955: the counted text invokes that named definition of the same article as an input, which is not retained in this card, so the reference is declared as a pointer carrying the number the article itself prints for it (the section and the place of the statement inside that section, derived by the census v3.5 rule for the shared section-scoped theorem counter); the cited definition is an allowed input and none of it is reproduced inside the span.\par\medskip
\ifdefined\hypertarget\hypertarget{sourcepointer-6}{}\fi
\noindent\textit{Source reference.} Original equation 4.5: operations/codex-takeover-01a12558/corpus-repair-20261010/research-line/harness-ready-20261011/sources/sigma-2025-030-sign-convention-for-a-infty-operations-in-bo/pinned-original-source.tex, source lines 850--850. Navigation pointer invoked at line(s) 965: the counted text cites that numbered equation of the same article, which is not retained in this card, so the reference is declared as a pointer carrying the number the article itself prints for it (the section and the place of the equation inside that section, derived by the census v3.5 rule for the section-scoped equation counter); the cited equation is an allowed input and none of it is reproduced inside the span.\par\medskip
\begin{thm}\label{conclusion}
The operations ${\mathfrak m}_k$, $k = 0, 1, \dots$, that is the Bott--Morse $A_{\infty}$-operation in the de Rham model, in Definition $\ref{def}$ satisfy the filtered $A_{\infty}$-relation
\begin{equation*}
\sum_{k' + k'' = k+1} {\mathfrak m}_{k'} \circ \widehat{\mathfrak m}_{k''} (\xi_1, \dots, \xi_k) = 0 \qquad {\rm for} \ k=1,2, \dots.
\end{equation*}
\end{thm}

\begin{proof}
By Proposition \ref{bdryori}, we find the following.

\begin{clm}\label{rel3}
The summation of the right hand side of~\eqref{lem} over $k'$, $k''$, $B'$, $B''$ such that
${k' + k'' = k+1}$, $B=B' + B''$, $(k', B'), (k'', B'') \neq (1,0)$ is equal to
\[ (-1)^{\kappa + \kappa'} \bigl({\rm ev}^B_0\vert_{\partial {\mathcal M}_{k+1}(B;{\mathcal L}; {\mathcal R})}\bigr)_! \bigl({\rm ev}_1^{B*}\xi_1 \wedge \dots\wedge {\rm ev}_k^{B*} \xi_k\bigr).\]
\end{clm}

Note that
\begin{align*}
\kappa + \kappa' &{} \equiv \epsilon(\xi_1 ,\dots, \xi_k) + 1 + k + \sum_{i=1}^k \deg \xi_i + \dim R_{\alpha_0} + \mu(R_{\alpha_0}) - \sum_{p=1}^k
\mu\bigl(R_{\alpha_p}\bigr)  \\
&{} \equiv \epsilon(\xi_1 ,\dots, \xi_k) + 1 + \dim {\mathcal M}_{k+1}(B;{\mathcal L};{\mathcal R}) + \sum_{i=1}^k \deg \xi_i \mod 2 .
\end{align*}

Using Proposition \ref{corrStokes}, we have
\begin{gather}
{\rm d} \bigl({\rm ev}_0^B\bigr)_!\bigl({\rm ev}_1^{B*}\xi_1 \wedge \dots \wedge {\rm ev}_k^{B*}\xi_k\bigr)  =  \bigl({\rm ev}_0^B\bigr)_! {\rm d}\bigl({\rm ev}_1^{B*}\xi_1 \wedge \dots \wedge {\rm ev}_k^{B*} \xi_k\bigr) \nonumber \\
\qquad{} + (-1)^\nu \bigl({\rm ev}_0^B\vert_{\partial {\mathcal M}_{k+1}(B;{\mathcal L};{\mathcal R})}\bigr)_!
\bigl({\rm ev}_1^{B*} \xi_1 \wedge \dots \wedge {\rm ev}_k^{B*}\xi_k\bigr),\label{Stokes}
\end{gather}
where $\nu = \dim {\mathcal M}_{k+1}(B;{\mathcal L};{\mathcal R}) + \sum_{i=1}^k \deg \xi_i.$

Combining~\eqref{rel1},~\eqref{rel2}, Claim \ref{rel3} and~\eqref{Stokes}, we have
\begin{gather*}
 {\mathfrak m}_{1,0} \circ {\mathfrak m}_{k, B} (\xi_1, \dots, \xi_k) + {\mathfrak m}_{k,B} \circ \hat{\mathfrak m}_{1,0} (\xi_1, \dots, \xi_k)  \\
\qquad{} + \sum_{(k',B'), (k'', B'') \neq (1,0)} {\mathfrak m}_{k',B'} \circ \hat{\mathfrak m}_{k'',B''} (\xi_1, \dots, \xi_k) = 0
\end{gather*}
for all $(k,B) \neq (1,0)$.
Recall that, in the case that $(k, B) = (1,0)$, ${\mathfrak m}_{1,0}= {\rm d}$ clearly satisfies ${\mathfrak m}_{1,0} \circ {\mathfrak m}_{1,0} = 0$.
Hence, we obtain Theorem \ref{conclusion}.
\end{proof}

\end{document}
